Theorem 2. In Theorem 1, if $m=n$, then there is a homeomorphism $g: \mathbf{R}^{n} \leftrightarrow \mathbf{R}^{n}$ such that $g \mid \boldsymbol{\sigma}^{n}$ is a simplicial homeomorphism $\boldsymbol{\sigma}^{n} \leftrightarrow \tau^{n}$.

Proof. The mapping $v \mapsto v-v_{0}$ is a homeomorphism $\mathbf{R}^{n} \leftrightarrow \mathbf{R}^{n}$, and maps every simplex simplicially onto a simplex. The composition of two such mappings has the same properties. Therefore we may assume, with no loss of generality, that $v_{0}$ is the origin in $\mathbf{R}^{n}$. Similarly for $w_{0}$. It follows that $\left\{v_{1}, v_{2}, \ldots, v_{n}\right\}$ and $\left\{w_{1}, w_{2}, \ldots, w_{n}\right\}$ are linearly independent. Now for every

$$
v=\sum_{i=1}^{n} \alpha_{i} v_{i} \in \mathbf{R}^{n}
$$

3 The Schönflies theorem for polygons in $\mathbf{R}^{2}$\\
we define

$$
g(v)=\sum_{i=1}^{n} \alpha_{i} w_{i} .
$$

Now $g \mid \sigma^{n}$ is the $f$ of Theorem 1. \(\square\)
